\documentclass[11pt]{article}
\usepackage[utf8]{inputenc}
\usepackage{geometry}
\geometry{a4paper, margin=1in}
\usepackage{hyperref}
\usepackage{booktabs}
\usepackage{enumitem}
\usepackage{amsmath}
\usepackage{float}
\usepackage{xcolor}

\title{Lifecycle 1 --- Q\&A Preparation \\ \large Questions Your Guide Will Ask}
\author{Study Reference}
\date{2026-03-18}

\begin{document}

\maketitle

\tableofcontents
\newpage

% ============================================================================
\section{Basic Concepts (Foundation Questions)}
% ============================================================================

\subsection{Q: What is a Convolutional Neural Network (CNN)?}

A CNN is a neural network designed for images. It slides small filters (kernels) across the image to detect patterns:

\begin{itemize}[noitemsep]
    \item \textbf{Early layers} detect simple features: edges, corners, color gradients
    \item \textbf{Middle layers} detect textures: soil patterns, vegetation edges
    \item \textbf{Deep layers} detect complex objects: landslide scars, debris trails
\end{itemize}

\textbf{Key operation --- Convolution:} A 3$\times$3 filter slides across the image, multiplying pixel values by filter weights and summing them. This produces a ``feature map'' highlighting where that pattern appears.

\textbf{Why CNNs for images:} They exploit \textit{spatial locality} (nearby pixels are related) and \textit{weight sharing} (same filter applied everywhere), making them efficient for image data.

\subsection{Q: What is backpropagation?}

Backpropagation is how the network \textit{learns}:

\begin{enumerate}[noitemsep]
    \item \textbf{Forward pass:} Image goes through network, produces prediction
    \item \textbf{Loss calculation:} Compare prediction with true label (Cross-Entropy Loss)
    \item \textbf{Backward pass:} Calculate how much each weight contributed to the error (chain rule)
    \item \textbf{Weight update:} Adjust weights to reduce error (optimizer like Adam/SGD)
\end{enumerate}

\textbf{Analogy:} A student takes an exam (forward pass), gets a score (loss), sees which answers were wrong (backward pass), and studies those topics more (weight update).

\subsection{Q: What is softmax?}

Softmax converts raw network outputs (logits) into probabilities that sum to 1:
\[
\text{softmax}(z_i) = \frac{e^{z_i}}{\sum_j e^{z_j}}
\]

\textbf{Example:} Logits $[2.0, 0.5]$ for [non-landslide, landslide]:
\[
P(\text{non-landslide}) = \frac{e^{2.0}}{e^{2.0} + e^{0.5}} = \frac{7.39}{9.04} = 81.7\%
\]

The class with highest probability is the prediction.

\subsection{Q: What is Cross-Entropy Loss?}

Measures how wrong the prediction is:
\[
\text{CE} = -\log(p_{\text{true class}})
\]

\begin{itemize}[noitemsep]
    \item Model predicts 90\% for correct class: loss = $-\log(0.9) = 0.105$ (low, good)
    \item Model predicts 10\% for correct class: loss = $-\log(0.1) = 2.302$ (high, bad)
\end{itemize}

The network minimizes this during training.

\subsection{Q: What is the difference between accuracy, precision, recall, F1?}

\begin{table}[H]
    \centering
    \renewcommand{\arraystretch}{1.3}
    \begin{tabular}{lp{10cm}}
        \toprule
        \textbf{Metric} & \textbf{What it measures} \\
        \midrule
        \textbf{Accuracy} & Out of ALL images, how many did we classify correctly? \\
        \textbf{Precision} & Out of all images we \textit{flagged as landslide}, how many actually were? (Low = too many false alarms) \\
        \textbf{Recall} & Out of all \textit{actual landslides}, how many did we catch? (Low = missed landslides = \textbf{dangerous}) \\
        \textbf{F1-Score} & Harmonic mean of precision and recall. Balances both. \\
        \bottomrule
    \end{tabular}
\end{table}

\textbf{Why recall matters most:} A missed landslide can cost lives. A false alarm only wastes resources.

\subsection{Q: What is ROC-AUC?}

\begin{itemize}[noitemsep]
    \item \textbf{ROC curve:} Plot of True Positive Rate vs False Positive Rate at every threshold
    \item \textbf{AUC:} Area Under this Curve (0.5 = random, 1.0 = perfect)
    \item \textbf{Meaning:} AUC = 85.53\% means the model ranks a random landslide image higher than a random non-landslide image 85.53\% of the time
\end{itemize}

\textbf{Why not just accuracy?} With 72\% non-landslide images, always predicting ``no landslide'' gives 72\% accuracy but catches zero landslides. AUC and F1 are not fooled by this.

% ============================================================================
\section{Architecture Questions}
% ============================================================================

\subsection{Q: Why did you start with AlexNet?}

\begin{enumerate}[noitemsep]
    \item \textbf{Baseline:} Need a performance floor. Every improvement must beat this.
    \item \textbf{Simple and fast:} Trains quickly, lets us verify the entire pipeline works end-to-end
    \item \textbf{Historical significance:} AlexNet (2012) started the deep learning revolution
\end{enumerate}

\subsection{Q: What is transfer learning? Why does it help?}

Using a model pre-trained on ImageNet (1.2M images) and fine-tuning on our 2,190 images.

\textbf{Why:} ImageNet teaches edges, textures, shapes. Our small dataset only needs to teach landslide-specific patterns. Without it, the model must learn everything from too few images.

\textbf{Analogy:} Hiring an experienced photographer and teaching them landslides, vs training someone who has never seen a photo.

\subsection{Q: What is EfficientNet's compound scaling?}

Traditional models scale only depth OR width OR resolution. EfficientNet scales all three:
\begin{itemize}[noitemsep]
    \item Depth $= \alpha^\phi$ (more layers)
    \item Width $= \beta^\phi$ (more filters per layer)
    \item Resolution $= \gamma^\phi$ (larger input images)
\end{itemize}

Result: EfficientNet-B3 achieves better accuracy than AlexNet with \textbf{5$\times$ fewer parameters} (12M vs 62M).

\subsection{Q: Why EfficientNet-B3 accuracy (79.97\%) is lower than pretrained AlexNet (80.11\%) but you say it's better?}

Because accuracy alone is misleading. Look at the full picture:
\begin{itemize}[noitemsep]
    \item \textbf{Recall:} 78.20\% vs 67.30\% --- EfficientNet catches \textbf{11\% more landslides}
    \item \textbf{F1-Score:} 70.21\% vs 66.98\% --- better precision-recall balance
    \item \textbf{ROC-AUC:} 88.93\% vs 85.73\% --- much better ranking ability
\end{itemize}

AlexNet pretrained had higher accuracy by being overly conservative (predicting non-landslide more often). EfficientNet is more aggressive at detecting landslides, which is what we want.

% ============================================================================
\section{HMM Questions}
% ============================================================================

\subsection{Q: What is an HMM? How does it work?}

A Hidden Markov Model models sequences where the true state is ``hidden'':

\begin{itemize}[noitemsep]
    \item \textbf{Hidden states:} True landslide regime (Shallow Slide, Deep Slide, Debris Flow, Rockfall) --- can't observe directly
    \item \textbf{Observations:} What we see --- landslide type recorded in NASA catalog
    \item \textbf{Transition matrix:} Probability of moving between states (e.g., after shallow slide, 30\% chance of debris flow)
    \item \textbf{Emission matrix:} Probability of observing a type given the hidden state
\end{itemize}

\subsection{Q: What is the Baum-Welch algorithm?}

An iterative training algorithm for HMMs (like EM algorithm):
\begin{enumerate}[noitemsep]
    \item \textbf{E-step:} Estimate which hidden states are most likely given current parameters
    \item \textbf{M-step:} Update transition/emission matrices to best explain the estimated states
    \item Repeat until convergence (log-likelihood stops improving)
\end{enumerate}

We ran 100 iterations, reaching log-likelihood of $-7314.96$.

\subsection{Q: Why did you choose 4 hidden states for the HMM?}

Four states map to the major landslide regimes in geotechnical literature: Shallow Landslide, Deep Landslide, Debris Flow, Rockfall. This was later increased to 8 states in LC2 to capture finer distinctions (monsoon-triggered vs earthquake-triggered events).

% ============================================================================
\section{Dataset Questions}
% ============================================================================

\subsection{Q: Isn't 2,190 images too small for deep learning?}

Yes, which is exactly why we use transfer learning (ImageNet pre-training provides a foundation). The small dataset size motivated several design choices:
\begin{itemize}[noitemsep]
    \item Transfer learning instead of training from scratch
    \item Data augmentation during training (random flips, rotations, color jitter)
    \item WeightedRandomSampler to handle class imbalance
    \item Ensemble methods to maximize performance from limited data
\end{itemize}

\subsection{Q: Why binary classification instead of multi-class?}

Binary (landslide vs non-landslide) provides a clean signal for measuring detection capability. Multi-class was attempted in LC3 and dropped to 71.67\% accuracy due to insufficient per-class samples. Binary classification is the pragmatic first step.

\subsection{Q: How did you split the data?}

The HR-GLDD dataset comes pre-split: 2,190 train / 550 validation / 699 test. We did NOT create our own splits --- this ensures reproducibility and fair comparison with other papers using the same dataset.

\end{document}
